Sea $\varphi:D\subseteq\mathbb R^n\to\mathbb R$ de clase $C^1$ y sea $\mathbf{F}=\nabla\varphi$. Suponga que una curva regular por partes $C$ está contenida en el conjunto de nivel $\varphi(\mathbf{x})=c$.
\begin{enumerate}
    \item Demuestre mediante el teorema fundamental de las integrales de línea que
    \[
    \int_C\mathbf{F}\cdot d\vec r=0.
    \]
    \item Demuestre el mismo resultado directamente: para una parametrización $\vec r(t)$ de cada tramo, derive la identidad $\varphi(\vec r(t))=c$ y relacione el resultado con $\mathbf{F}(\vec r(t))\cdot\vec r\,'(t)$.
    \item Aplique ambas interpretaciones a $\mathbf{F}(x,y,z)=\langle2x,2y,2z\rangle$ y a cualquier curva contenida en la esfera $x^2+y^2+z^2=R^2$. Explique la relación geométrica entre el campo y el vector tangente a la curva.
\end{enumerate}
